\subsection{Continuum constraints through a Gaussian coefficient}\label{sec:methods}

To express a continuum restriction through one coefficient, we first
choose a function in the retained span whose selected derivative,
including the function itself at order zero, is positive throughout $I$.
A sufficiently large coefficient along this function then makes the
selected quantity nonnegative everywhere on the interval, whatever
the values of the remaining coefficients. Since an invertible change
of coordinates preserves every admissible coefficient vector in the
finite space, the restriction can be represented by a single lower
bound without changing the underlying Gaussian law. The following
result integrates that coefficient and reconstructs it conditionally,
using the standard normal distribution function $\Phi$ and survival
function $\Phi^{\mathrm c}(x)=1-\Phi(x)$.

\begin{theoremEnd}[combined main,category=finite]{theorem}[Gaussian conditioning under a continuum sign restriction]
\CombinedResultCheck{thm:finite}{theorem}\label{thm:finite}
Under the finite Gaussian model in \cref{sec:finite-model}, let
$\boldsymbol a=(a_0,\ldots,a_m)^T$ denote the coefficient vector and fix
$r\in\{0,1,2\}$. Suppose an invertible linear change of coordinates
$\boldsymbol a\mapsto(\boldsymbol g,\gamma)\in\R^m\times\R$
gives $f_{\boldsymbol a}=f_{\boldsymbol g}+\gamma h_r$.
Here $h_r$ is a fixed function in the retained span, chosen to have a
uniformly positive $r$th derivative on $I$, and $\gamma$ is its scalar
coefficient. The curve $f_{\boldsymbol g}$ depends linearly on the
remaining coordinates $\boldsymbol g$. Assume these functions belong
to $C^r(I)$ with
$\min_{x\in I}\frac{\mathrm{d}^rh_r}{\mathrm{d}x^r}(x)>0$.
In these coordinates, the ordinary Gaussian posterior has the form
\begin{equation}\label{eq:conditional}
 \boldsymbol g\mid\boldsymbol y\sim N(\boldsymbol m_g,\boldsymbol V_g),
 \qquad \gamma\mid\boldsymbol g,\boldsymbol y\sim N(\mu(\boldsymbol g),s^2).
\end{equation}
Here $\boldsymbol V_g$ is positive definite, $\mu$ is affine, and
$s>0$ is constant in $\boldsymbol g$. The finite, attained threshold
\begin{equation}\label{eq:threshold}
 \tau_r(\boldsymbol g)=\max_{x\in I}
 \frac{-\frac{\mathrm{d}^rf_{\boldsymbol g}}{\mathrm{d}x^r}(x)}
      {\frac{\mathrm{d}^rh_r}{\mathrm{d}x^r}(x)}
\end{equation}
satisfies $f_{\boldsymbol a}\in\mathcal C_r$ exactly when
$\gamma\geq\tau_r(\boldsymbol g)$. This event has positive ordinary
posterior probability and a null boundary in coefficient space.
Its marginal density for $\boldsymbol g$ with
respect to Lebesgue measure on $\R^m$ is
\begin{equation}\label{eq:marginal}
 p(\boldsymbol g\mid\boldsymbol y,\mathcal C_r)=
 \frac{p(\boldsymbol g\mid\boldsymbol y)
 \Phi^{\mathrm c}((\tau_r(\boldsymbol g)-\mu(\boldsymbol g))/s)}
 {P(f_{\boldsymbol a}\in\mathcal C_r\mid\boldsymbol y)}.
\end{equation}
Given this draw of $\boldsymbol g$, draw $\gamma$ from the conditional normal in
\cref{eq:conditional} truncated below at $\tau_r(\boldsymbol g)$ and
apply the inverse coefficient map. Equivalently, for an independent
$U\sim\operatorname{Uniform}(0,1)$,
\begin{equation}\label{eq:reconstruction}
 \gamma=\mu(\boldsymbol g)-s\Phi^{-1}
 \bigl(U\Phi^{\mathrm c}((\tau_r(\boldsymbol g)-\mu(\boldsymbol g))/s)\bigr).
\end{equation}
This reconstructs the full conditioned coefficient law, under which
$\min_I\frac{\mathrm{d}^rf_{\boldsymbol a}}{\mathrm{d}x^r}>0$ almost surely.
\end{theoremEnd}

\begin{proofEnd}
\CombinedProofCheck{thm:finite}
Fix $r$, $I$, the invertible map, and the ordinary posterior under
the finite Gaussian model in \cref{sec:finite-model} and the
positive-direction hypotheses of
\cref{thm:finite}, and write
$b_r=\min_{x\in I}\frac{\mathrm{d}^rh_r}{\mathrm{d}x^r}(x)>0$.
The role of the positive direction is to turn every
derivative inequality into a lower bound on the same scalar coefficient:
division by $\frac{\mathrm{d}^rh_r}{\mathrm{d}x^r}$ gives the threshold equivalence,
and continuity on compact $I$, together with $b_r>0$, makes
\cref{eq:threshold} finite and attained. Since the invertible map
preserves positive definiteness of the posterior covariance, Gaussian conditioning
\citep[Section 2.2]{rasmussen2006gpml} gives
\cref{eq:conditional}, including a positive Schur complement $s^2$.

Integrating the restricted joint Gaussian density over the scalar
coefficient now leaves its conditional survival probability as a
weight on $\boldsymbol g$, which proves \cref{eq:marginal}.
This weight is strictly positive for every $\boldsymbol g$, so the
conditioning event has positive probability. Given a draw from the
weighted marginal, conditioning the Gaussian law of
$\gamma\mid\boldsymbol g,\boldsymbol y$ on
$\gamma\geq\tau_r(\boldsymbol g)$ gives the stated truncated normal,
whose survival-function inverse is \cref{eq:reconstruction}.
Applying the inverse coefficient map therefore recovers the original
coefficient law conditional on the sign restriction.

The minimum derivative is continuous in the coefficients, so coefficient
vectors with a positive minimum lie in the event's interior and those
with a negative minimum lie in its exterior. Conditional on
$\boldsymbol g$, the nondegenerate normal assigns zero
mass to equality at the threshold. At equality an attained maximizing
location has zero derivative, whereas above it
\[
 \min_I\frac{\mathrm{d}^rf_{\boldsymbol a}}{\mathrm{d}x^r}
 \geq(\gamma-\tau_r(\boldsymbol g))b_r>0.
\]
At equality, an arbitrarily small increase or decrease in $\gamma$
moves to either side of the event, identifying its coefficient boundary
as precisely the threshold graph.
Its zero Gaussian probability implies that the reconstructed law lies
strictly above the graph almost surely, proving the strict-derivative
assertion.
\end{proofEnd}


The Gaussian tail in \cref{eq:marginal} is the conditional probability
that the separated coefficient is large enough to satisfy the
constraint, so values of $\boldsymbol g$ receive less posterior weight
when they require an unlikely value of that coefficient. Conditional
reconstruction retains the resulting dependence on $\boldsymbol g$ and
therefore gives joint draws of the curve and its derivatives. This
dependence is inherited from the original Gaussian model and must be
transported through the coefficient map, since assigning a new
independent prior to the transformed coordinates would change the law.

A related conditioning structure appears in GP regression with censored
observations, where a Gaussian posterior density is weighted by the
conditional probability of the censoring event
\citep{gormhoffmann2025censored}. In that setting, inequalities describe
censored measurements at finitely many locations, whereas the threshold
here represents a restriction on the latent curve throughout an interval.

The null-boundary statement in \cref{thm:finite} means that, under
its assumptions, the conditioned law assigns probability zero to curves
whose minimum response or selected derivative equals zero. Thus a draw
from this law has a strictly positive minimum almost surely, although
these minima can be arbitrarily close to zero. For monotonicity,
positivity of the derivative at every point is stronger than strict
ordering of function values, which alone does not require a positive
derivative everywhere. Nearly flat regions therefore remain possible,
but the analyticity of the finite expansion means that assigning
positive probability to exact plateaus or linear segments of positive
length requires a different prior.

The random intercept supplies a constant direction for positivity,
whereas monotonicity and convexity require directions with positive
slope and curvature. For each choice, evaluating the threshold still
requires finding a global extremum over $I$. On the padded interval of
length $2L$, the zero boundary conditions require
$\sin(2L\omega_j)=0$, giving
$\omega_j=j\pi/(2L)=j\omega_1$ for $j\geq1$. The basis and its
derivatives therefore involve sine and cosine harmonics of a common
fundamental frequency.
We introduce
\[
 \theta(x)=\omega_1(x-c+L),\qquad w(x)=\cos\theta(x),
\]
because each cosine harmonic becomes a polynomial in $w$, while each
sine harmonic becomes $\sin\theta$ times a polynomial in $w$, through
identities for the Chebyshev polynomials in \cref{def:chebyshev}.
These identities allow the stationary locations needed for the
thresholds to be characterized by polynomial equations.

Since $I$ lies strictly inside the padded domain, $0<\theta(x)<\pi$
and $\mathrm{d}w/\mathrm{d}x=-\omega_1\sin\theta(x)<0$ throughout $I$. Thus $w$ maps
$I$ one-to-one onto a compact interval in $(-1,1)$, reversing the order
of its endpoints. Its nonzero derivative ensures that interior
stationary locations correspond in the two coordinates, so the change
of variable preserves the complete search for a continuum extremum.

\begin{definition}[Chebyshev polynomials]\label{def:chebyshev}
The Chebyshev polynomials of the first kind, $T_j$, and second kind,
$U_j$, are defined for $w\in\R$ by the initial values and recurrences
\[
 \begin{aligned}
 T_0(w)&=1,\quad T_1(w)=w,
 & T_{j+1}(w)&=2wT_j(w)-T_{j-1}(w),\\
 U_0(w)&=1,\quad U_1(w)=2w,
 & U_{j+1}(w)&=2wU_j(w)-U_{j-1}(w),
 \end{aligned}
\]
where the recurrences hold for integers $j\geq1$.
\end{definition}

For $w=\cos\theta$, these polynomials satisfy the standard trigonometric
identities \citep[Equations (18.5.1) and (18.5.2)]{nist2026dlmf}
\[
 T_j(w)=\cos(j\theta),\qquad
 \sin\theta\,U_{j-1}(w)=\sin(j\theta),\qquad j\geq1.
\]
In particular, the basis and its first two derivatives satisfy
\[
 \phi_{j,L}(x)=\frac{\sin\theta(x)}{\sqrt L}\,U_{j-1}(w(x)),\quad
 \frac{\mathrm{d}\phi_{j,L}}{\mathrm{d}x}(x)=\frac{j\omega_1}{\sqrt L}\,T_j(w(x)),\quad
 \frac{\mathrm{d}^2\phi_{j,L}}{\mathrm{d}x^2}(x)=-\omega_j^2\phi_{j,L}(x).
\]
The second derivative therefore shares the sine factor in the first
identity. For positivity, the polynomial
representing the first derivative identifies the stationary locations
of the curve. With the directions chosen below, monotonicity gives a
ratio of polynomials with a positive denominator, while convexity gives
a polynomial after the common sine factor cancels. Differentiating
these expressions reduces their stationary-point conditions to
polynomial equations. Evaluating the interval
endpoints and these locations determines the global thresholds, with
constant cases handled explicitly in the corollaries.

The following three corollaries apply \cref{thm:finite} to positivity,
monotonicity, and convexity, specifying a suitable direction and the
corresponding continuum threshold for each restriction.

\begin{theoremEnd}[combined main,category=positive]{corollary}[Positivity in the Laplacian expansion]
\CombinedResultCheck{cor:positive}{corollary}\label{cor:positive}
Under the finite Gaussian model in \cref{sec:finite-model}, suppose
$m\geq1$ and $I\subset(c-L,c+L)$. Set
$h_0=1$, $\boldsymbol g=(a_1,\ldots,a_m)^T$, and
$f_{\boldsymbol g}=\sum_{j=1}^ma_j\phi_{j,L}$.
Then $f_{\boldsymbol a}=f_{\boldsymbol g}+a_0h_0$, with inverse
obtained by reinserting $a_0$.

For $\theta(x)=\omega_1(x-c+L)$, $w(x)=\cos\theta(x)$, and
$J=[w(x_+),w(x_-)]\subset(-1,1)$, define
$P_{0,\boldsymbol g}(w)=\sum_{j=1}^mja_jT_j(w)$. The threshold is
\begin{equation}\label{eq:positive-threshold}
 \tau_0(\boldsymbol g)=-\min_{x\in I}f_{\boldsymbol g}(x).
\end{equation}
\cref{thm:finite} applies with scalar coefficient $\gamma=a_0$.

The threshold is attained at a physical endpoint or a location whose
cosine coordinate is an interior real root in $J$ of
$P_{0,\boldsymbol g}$. This stationary polynomial has degree at most
$m$ when nonzero. All root multiplicities and reduced degrees are
included. If the stationary polynomial vanishes identically, the
curve $f_{\boldsymbol g}$ is constant and either endpoint suffices.
\end{theoremEnd}

\begin{proofEnd}
\CombinedProofCheck{cor:positive}
Use the finite Gaussian model in \cref{sec:finite-model} and its compact
physical interval. Separating the intercept is an invertible permutation
of the coefficients, with $\min_Ih_0=1$, so \cref{thm:finite} at
order zero gives \cref{eq:positive-threshold} and the conditioned law.
To locate the minimum over the whole interval, use the Chebyshev
identities to write
\[
 f_{\boldsymbol g}(x)=\frac{\sqrt{1-w(x)^2}}{\sqrt L}
             \sum_{j=1}^ma_jU_{j-1}(w(x)),
 \qquad
 \frac{\mathrm{d}f_{\boldsymbol g}}{\mathrm{d}x}(x)=
             \frac{\omega_1}{\sqrt L}P_{0,\boldsymbol g}(w(x)).
\]
Because the cosine is strictly decreasing for $0<\theta(x)<\pi$,
interior stationary locations correspond exactly to interior roots of
the displayed polynomial. The extreme value theorem and Fermat's
theorem therefore give the endpoint and stationary candidates, with
finitely many distinct roots, including repeated roots, when the
polynomial is nonzero and has degree at most $m$. If it vanishes
identically, the derivative is zero throughout $I$, so the mean value
theorem gives constancy. Although the intercept is
independent of the spectral coefficients in the prior, reconstruction
uses its conditional posterior given $\boldsymbol g$ and the data.
\end{proofEnd}

\begin{theoremEnd}[combined main,category=monotone]{corollary}[Monotonicity in the Laplacian expansion]
\CombinedResultCheck{cor:monotone}{corollary}\label{cor:monotone}
Under the finite Gaussian model in \cref{sec:finite-model}, suppose
$m\geq4$ and $I\subset[c-L/2,c+L/2]$. Set
$h_1=-\phi_{2,L}-\phi_{4,L}/4$,
$\boldsymbol g=(g_0,g_1,g_3,\ldots,g_m)^T$, and
$f_{\boldsymbol g}=g_0+\sum_{j=1}^mg_j\phi_{j,L}$, where
$g_2=0$, $g_4=a_4-a_2/4$, and $g_j=a_j$ for $j\notin\{2,4\}$.
Then $f_{\boldsymbol a}=f_{\boldsymbol g}+(-a_2)h_1$, with inverse
obtained by reinserting $a_2$ and setting $a_4=g_4+a_2/4$.

For $\theta(x)=\omega_1(x-c+L)$, $w(x)=\cos\theta(x)$, and
$J=[w(x_+),w(x_-)]\subset[-1/\sqrt2,1/\sqrt2]$, define
\begin{equation}\label{eq:monotone-polynomials}
 P_{\boldsymbol g}(w)=\sum_{\substack{1\leq j\leq m\\j\ne2}}jg_jT_j(w),
 \qquad D(w)=1+4w^2-8w^4.
\end{equation}
Since $D\geq1$ on $J$, the threshold is
\begin{equation}\label{eq:monotone-threshold}
 \tau_1(\boldsymbol g)=\max_{w\in J}\frac{-P_{\boldsymbol g}(w)}{D(w)}.
\end{equation}
\cref{thm:finite} applies with scalar coefficient $\gamma=-a_2$.

The threshold is attained at a physical endpoint or a location whose
cosine coordinate is an interior real root in $J$ of
$\frac{\mathrm{d}P_{\boldsymbol g}}{\mathrm{d}w}D-P_{\boldsymbol g}\frac{\mathrm{d}D}{\mathrm{d}w}$.
This stationary polynomial has degree at most $m+3$ when nonzero.
All root multiplicities and reduced degrees are included. If the
stationary polynomial vanishes identically, the quotient
$-P_{\boldsymbol g}/D$ is constant and either endpoint suffices.
A nonzero constant stationary polynomial gives only endpoint candidates.
\end{theoremEnd}

\begin{proofEnd}
\CombinedProofCheck{cor:monotone}
Use the finite Gaussian model in \cref{sec:finite-model} with the rank and closed padding interval in
\cref{cor:monotone}. Substitution of the inverse map proves the curve
identity and invertibility, while the identity
$T_j(w(x))=\cos(j\theta(x))$ gives the derivative formulas
\[
 \frac{\mathrm{d}f_{\boldsymbol g}}{\mathrm{d}x}(x)=\frac{\omega_1}{\sqrt L}P_{\boldsymbol g}(w(x)),
 \qquad \frac{\mathrm{d}h_1}{\mathrm{d}x}(x)=\frac{\omega_1}{\sqrt L}D(w(x)).
\]
The assumed padding confines $\theta(x)$ to $[\pi/4,3\pi/4]$, where
the cosine is strictly decreasing and
$D(w)=1+4w^2(1-2w^2)\geq1$, including at the physical endpoints.
The positive derivative bound therefore allows \cref{thm:finite}
to supply the conditioned law and strictness.
Because the map transports the original Gaussian covariance,
$g_4$ and $a_2$ retain their induced dependence.

Since $D>0$ on compact $J$, the extreme value theorem and Fermat's
theorem locate the maximum of $-P_{\boldsymbol g}/D$ at an endpoint
or an interior stationary point. The quotient rule gives the
stationary polynomial in the statement, whose degree is at most
$m+3$ because $P_{\boldsymbol g}$ and $D$ have degrees at most
$m$ and $4$. If this polynomial vanishes identically, the quotient
is constant by the mean value theorem and either endpoint suffices.
Otherwise it has finitely many distinct roots, each a candidate
when it lies in the interior of $J$, regardless of multiplicity.
This also covers reduced degrees, including a nonzero constant
stationary polynomial with no interior candidates.
\end{proofEnd}

\begin{theoremEnd}[combined main,category=convex]{corollary}[Convexity in the Laplacian expansion]
\CombinedResultCheck{cor:convex}{corollary}\label{cor:convex}
Under the finite Gaussian model in \cref{sec:finite-model}, suppose
$m\geq2$ and $I\subset(c-L,c+L)$. Set $h_2=-\phi_{1,L}$,
$\boldsymbol g=(a_0,a_2,\ldots,a_m)^T$, and
$f_{\boldsymbol g}=a_0+\sum_{j=2}^ma_j\phi_{j,L}$.
Then $f_{\boldsymbol a}=f_{\boldsymbol g}+(-a_1)h_2$, with inverse
obtained by reinserting $a_1$.

For $\theta(x)=\omega_1(x-c+L)$, $w(x)=\cos\theta(x)$, and
$J=[w(x_+),w(x_-)]\subset(-1,1)$, define the polynomial and threshold
\begin{equation}\label{eq:convex-threshold}
 Q_{\boldsymbol g}(w)=\sum_{j=2}^mj^2a_jU_{j-1}(w),
 \qquad \tau_2(\boldsymbol g)=\max_{w\in J}Q_{\boldsymbol g}(w).
\end{equation}
\cref{thm:finite} applies with scalar coefficient $\gamma=-a_1$.

The threshold is attained at a physical endpoint or a location whose
cosine coordinate is an interior real root in $J$ of
$\frac{\mathrm{d}Q_{\boldsymbol g}}{\mathrm{d}w}$. The polynomial $Q_{\boldsymbol g}$
has degree at most $m-1$, so its stationary polynomial has degree at
most $m-2$ when nonzero. All root multiplicities and reduced degrees
are included. If the stationary polynomial vanishes identically,
$Q_{\boldsymbol g}$ is constant and either endpoint suffices.
A nonzero constant stationary polynomial gives only endpoint candidates.
\end{theoremEnd}

\begin{proofEnd}
\CombinedProofCheck{cor:convex}
Use the finite Gaussian model in \cref{sec:finite-model} with the rank and open padding interval in
\cref{cor:convex}. The coordinate map is an invertible signed
permutation. Since $0<\theta(x)<\pi$ on compact $I$,
\[
 \frac{\mathrm{d}^2h_2}{\mathrm{d}x^2}(x)=\frac{\omega_1^2}{\sqrt L}\sin\theta(x)>0,
 \qquad
 \frac{\mathrm{d}^2f_{\boldsymbol g}}{\mathrm{d}x^2}(x)=
 -\frac{\omega_1^2}{\sqrt L}\sin\theta(x)\,Q_{\boldsymbol g}(w(x)).
\]
The second formula uses
$\sin(j\theta(x))=\sin\theta(x)\,U_{j-1}(w(x))$, while compactness
keeps the positive curvature in the first formula bounded away from
zero. Thus \cref{thm:finite} gives the claimed threshold and posterior
under the stated open padding condition.

The extreme value theorem and Fermat's theorem locate the maximum
of $Q_{\boldsymbol g}$ at an endpoint or an interior zero of
$\frac{\mathrm{d}Q_{\boldsymbol g}}{\mathrm{d}w}$. The degree of $Q_{\boldsymbol g}$
is at most $m-1$, so this derivative has degree at most $m-2$ when
nonzero. If it vanishes identically, $Q_{\boldsymbol g}$ is constant
and either endpoint suffices. Otherwise its finitely many distinct
interior roots enter the candidate set regardless of multiplicity,
with none when the derivative is a nonzero constant. The same
argument covers every reduced degree.
\end{proofEnd}


For positivity, separating $a_0$ leaves the spectral curve
$f_{\boldsymbol g}$, and \cref{eq:positive-threshold} is the vertical
shift needed to make its minimum nonnegative. The intercept's
conditional posterior mean can depend on the spectral coefficients,
so reconstruction uses that conditional law even though the prior
coordinates are independent. The monotonicity map makes the same
need to preserve dependence explicit by replacing $a_4$ with
$g_4=a_4-a_2/4$, which is correlated with the separated coefficient
$-a_2$. The convexity map simply sets aside the first mode, but all
three constructions transport the ordinary Gaussian posterior before
conditioning. In the two derivative cases, the intercept stays among
the remaining coefficients and can change its posterior through
correlation, despite contributing neither slope nor curvature.

Although all three constructions use the same padded basis domain
$(c-L,c+L)$, their directions give different sufficient conditions on
the physical interval. The constant positivity direction requires no
additional padding, and $h_2$ has positive curvature throughout the
open basis domain, whereas the slope of the chosen $h_1$ is bounded
away from zero on the central half, including its endpoints. The
conditions in \cref{tab:cases} reflect these properties of the selected
directions, so imposing the central-half condition on convexity would
unnecessarily restrict that finite result.

\begin{table}[htbp]
\centering
\setlength{\tabcolsep}{4pt}
\begin{tabular}{llll}
\toprule
Quantity & Positivity & Monotonicity & Convexity\\
\midrule
Selected order & $r=0$ & $r=1$ & $r=2$\\
Separated coefficient & $a_0$ & $-a_2$ & $-a_1$\\
Direction & $1$ & $-\phi_{2,L}-\phi_{4,L}/4$ & $-\phi_{1,L}$\\
Finite interval & $I\subset(c-L,c+L)$ & $I\subset[c-L/2,c+L/2]$ & $I\subset(c-L,c+L)$\\
Rank condition & $m\geq1$ & $m\geq4$ & $m\geq2$\\
Threshold & $-\min_I f_{\boldsymbol g}$ & $\max_J(-P_{\boldsymbol g}/D)$ & $\max_J Q_{\boldsymbol g}$\\
Polynomial degree bound & $m$ & $m+3$ & $m-2$\\
Limit space & $C(I)$ & $C^1(I)$ & $C^2(I)$\\
Mat\'ern smoothness & $\nu>0$ & $\nu>1$ & $\nu>2$\\
\bottomrule
\end{tabular}
\caption{Positivity, monotonicity, and convexity as alternative
conditioning events in the same spectral model. Here $I$ is the physical
interval and $J$ its image under the decreasing cosine coordinate $w$.
The degree bounds apply to the nonzero polynomials whose roots identify
stationary locations.
The finite rank and
interval conditions are sufficient for the displayed directions. The
function-space limits additionally require \cref{ass:spectral} and the
padding and frequency-cutoff limits in \cref{thm:approximation}.}
\label{tab:cases}
\end{table}

To compute the thresholds in
\cref{eq:positive-threshold,eq:monotone-threshold,eq:convex-threshold},
we compare values at the two interval endpoints and at the interior
stationary locations identified by the polynomial roots in $J$.
For positivity, we evaluate $f_{\boldsymbol g}$ at the corresponding
locations in $I$ and take the negative of the smallest value. For
monotonicity and convexity, we take the largest value of
$-P_{\boldsymbol g}/D$ and $Q_{\boldsymbol g}$, respectively, at the
candidate values of $w$. Every interior real root must be included,
even when the derivative does not change sign there. If the stationary
polynomial vanishes identically, the function being optimized is
constant and an endpoint evaluation suffices.

Nonpositivity, nonincreasing regression, and concavity follow by negating
the observations and prior mean while retaining all covariances, then
applying the corresponding nonnegative-derivative construction and
negating the reconstructed curve and its derivatives. In particular,
concavity makes the slope nonincreasing but allows it to have either
sign, so a curvature restriction alone does not impose monotonicity.
\cref{tab:coefficient-reconstruction} summarizes the inverse maps for
all six individual sign constraints.

\begin{table}[htbp]
\centering
\small
\begin{tabular}{@{}lll@{}}
\toprule
Constraint & Reinserted coefficients & Remaining coefficients\\
\midrule
Positivity & $a_0=\gamma$
  & $a_j=g_j$, $1\leq j\leq m$\\
Negativity & $a_0=-\gamma$
  & $a_j=-g_j$, $1\leq j\leq m$\\
\addlinespace
Nondecreasing & $a_2=-\gamma$, $a_4=g_4-\gamma/4$
  & $a_j=g_j$, $j\notin\{2,4\}$\\
Nonincreasing & $a_2=\gamma$, $a_4=-g_4+\gamma/4$
  & $a_j=-g_j$, $j\notin\{2,4\}$\\
\addlinespace
Convexity & $a_1=-\gamma$
  & $a_j=g_j$, $j\ne1$\\
Concavity & $a_1=\gamma$
  & $a_j=-g_j$, $j\ne1$\\
\bottomrule
\end{tabular}
\caption[Reconstruction for six sign constraints]{Inverse maps to the
original coefficients for six constraints on $I$. Positivity and
negativity mean $f_{\boldsymbol a}\geq0$ and $f_{\boldsymbol a}\leq0$.
The negative, nonincreasing, and concave cases use
$(\boldsymbol g,\gamma)$ from the Gaussian calculation for
$-f_{\boldsymbol a}$ with negated observations and prior mean.
The other cases use $f_{\boldsymbol a}$. Indices range from $0$ to $m$
unless stated otherwise, and $g_j$ retains its original coefficient
index. Each pair uses the rank and padding conditions of its
corresponding case in \cref{cor:positive,cor:monotone,cor:convex}.}
\label{tab:coefficient-reconstruction}
\end{table}

Several restrictions can describe different aspects of the same
response. A nonnegative, nondecreasing, concave curve, for example,
has a nonnegative level and slope, with the slope nonincreasing across
the interval. We express such intersections by selecting derivative
orders and a sign at each order. For a nonempty index set
$A\subseteq\{0,1,2\}$ and a sign choice
$\varepsilon:A\to\{-1,1\}$, write $\varepsilon_r=\varepsilon(r)$ and define
\begin{equation}\label{eq:joint-event}
 \mathcal C_A^\varepsilon=\left\{f\in C^{\max A}(I):
       \varepsilon_r\frac{\mathrm{d}^rf}{\mathrm{d}x^r}(x)\geq0
       \text{ for every }r\in A\text{ and }x\in I\right\}.
\end{equation}
At order zero, the two signs give nonnegativity and nonpositivity.
At orders one and two, they give nondecreasingness and nonincreasingness,
and convexity and concavity, respectively. Taking $\varepsilon_r=1$
at every selected order gives positivity, nondecreasingness, and
convexity, their three pairwise combinations, and their triple
combination. Allowing either sign gives six individual cases and
20 joint combinations without opposing signs. Imposing both signs
at the same order instead forces the corresponding derivative to
vanish throughout $I$, an event of zero probability under the finite
Gaussian model.

To retain the scalar conditioning construction, we need one direction
$h$ in the retained span whose selected signed derivatives are all
uniformly positive on $I$. Increasing its coefficient then increases
every signed derivative involved in the restriction. Each order and
location imposes a lower bound on this same coefficient, so their
largest bound enforces the entire intersection. The direction
condition must be verified for the chosen combination, rank, and
interval. The following proposition gives the resulting Gaussian
conditioning law.

\begin{theoremEnd}[combined main,category=joint]{proposition}[Joint sign constraints]
\CombinedResultCheck{prop:joint}{proposition}\label{prop:joint}
Under the finite Gaussian model in \cref{sec:finite-model}, let
$\boldsymbol a=(a_0,\ldots,a_m)^T$ denote the coefficient vector and fix
a nonempty $A\subseteq\{0,1,2\}$ and signs
$\varepsilon:A\to\{-1,1\}$, with event $\mathcal C_A^\varepsilon$ defined in
\cref{eq:joint-event}. Suppose an invertible linear change of coordinates
$\boldsymbol a\mapsto(\boldsymbol g,\gamma)\in\R^m\times\R$
gives $f_{\boldsymbol a}=f_{\boldsymbol g}+\gamma h$.
Here $h$ is a fixed function in the retained span, shared by all selected
orders, $\gamma$ is its scalar coefficient, and $f_{\boldsymbol g}$
depends linearly on the remaining coordinates $\boldsymbol g$.
Assume $h,f_{\boldsymbol g}\in C^{\max A}(I)$ with
$\min_{x\in I}\varepsilon_r\frac{\mathrm{d}^rh}{\mathrm{d}x^r}(x)>0$
for every $r\in A$. Define
\begin{equation}\label{eq:joint-threshold}
 \eta_r(x;\boldsymbol g)=
       \frac{-\varepsilon_r\frac{\mathrm{d}^rf_{\boldsymbol g}}{\mathrm{d}x^r}(x)}
            {\varepsilon_r\frac{\mathrm{d}^rh}{\mathrm{d}x^r}(x)},
 \qquad
 \tau_A^\varepsilon(\boldsymbol g)=
       \max_{r\in A}\max_{x\in I}\eta_r(x;\boldsymbol g).
\end{equation}
This finite, attained threshold satisfies
$f_{\boldsymbol a}\in\mathcal C_A^\varepsilon$ exactly when
$\gamma\geq\tau_A^\varepsilon(\boldsymbol g)$.
The event has positive ordinary posterior probability and a null
coefficient boundary. Its marginal density and reconstruction are
\cref{eq:marginal,eq:reconstruction}, with $\tau_r,\mathcal C_r$
replaced by $\tau_A^\varepsilon,\mathcal C_A^\varepsilon$ and the
Gaussian parameters in \cref{eq:conditional} computed under this map.
The conditioned law satisfies
$\min_{x\in I}\varepsilon_r\frac{\mathrm{d}^rf_{\boldsymbol a}}{\mathrm{d}x^r}(x)>0$
almost surely for every $r\in A$. The envelope, strong log-concavity,
and conservative threshold-error conclusions extend under their
respective hypotheses.
\end{theoremEnd}

\begin{proofEnd}
\CombinedProofCheck{prop:joint}
Fix the coefficient map and signed direction satisfying the hypotheses
of \cref{prop:joint}.
Dividing each signed derivative inequality by its positive direction
derivative gives $\gamma\geq\eta_r(x;\boldsymbol g)$, so their intersection
gives \cref{eq:joint-threshold}. Compactness of $I$, finiteness
of $A$, and the denominator bounds make the maximum finite and attained.
The Gaussian conditioning argument in the proof of \cref{thm:finite} then gives
positive probability, the marginal density, and reconstruction.

Above the threshold, every selected signed derivative has a strictly
positive minimum. At equality, an attained maximizing pair gives a zero
derivative, and perturbing $\gamma$ crosses the event boundary. Thus the
boundary is the threshold graph, which has zero conditional Gaussian
probability for every $\boldsymbol g$, establishing almost-sure strictness. Each
$\eta_r(x;\boldsymbol g)$ is linear in $\boldsymbol g$, so the envelope
and strong log-concavity arguments apply with the compact index set
$A\times I$, and \cref{prop:error} applies to the same scalar truncation.
\end{proofEnd}


The joint restrictions can also reduce the number of continuum
searches. For the nonnegative, nondecreasing, concave example, take
$A=\{0,1,2\}$ with $\varepsilon_0=\varepsilon_1=1$ and
$\varepsilon_2=-1$. Concavity makes the slope smallest at $x_+$,
and nondecreasingness makes the response smallest at $x_-$.
Under the common map in \cref{prop:joint}, the threshold therefore becomes
\[
 \tau_A^\varepsilon(\boldsymbol g)
 =\max\left\{
       \max_{x\in I}\eta_2(x;\boldsymbol g),
       \eta_1(x_+;\boldsymbol g),\eta_0(x_-;\boldsymbol g)
       \right\}.
\]
Thus this combination requires one continuum curvature bound and
two endpoint bounds. All three use the same $h$ and $\boldsymbol g$,
whose choice determines the quotients in \cref{eq:joint-threshold}.
The supplementary material gives the HSGP coefficient map, conditions
on its direction coefficients, and the reductions for all sign choices,
with a table separating continuum and endpoint bounds for the 20
joint cases.
